% Build with slides/build.sh (upLaTeX + dvipdfmx).
% Theme: hkaomua/slide-template at 3d26f4f. FPGA evidence: 229aed0.
\documentclass[dvipdfmx,aspectratio=169,t]{beamer}
\usetheme{lineSeed}
\usepackage{colortbl,array,amsmath,pgfplots,listings}
\pgfplotsset{compat=1.18}
\input{metrics.tex}
\pdfstringdefDisableCommands{\def\translate#1{#1}}
% Keep the upstream heading rule visible above table/plot content.
\AddToHook{shipout/foreground}{%
  \ifnum\value{page}>1\relax
    \begin{tikzpicture}[remember picture,overlay]
      \fill[SeedRule,rounded corners=0.850395bp]
        ([xshift=36bp,yshift=-64.66732bp]current page.north west)
        rectangle ([xshift=684bp,yshift=-66.36811bp]current page.north west);
      \fill[SeedAccent,rounded corners=0.850395bp]
        ([xshift=36bp,yshift=-64.66732bp]current page.north west)
        rectangle ([xshift=87.02362bp,yshift=-66.36811bp]current page.north west);
    \end{tikzpicture}%
  \fi
}
\newcommand{\bigline}[1]{{\fontsize{24bp}{32bp}\selectfont\bfseries #1\par}}
\newcommand{\intent}[1]{\noindent{\bfseries #1\par}\vspace{8bp}}
\newcommand{\gap}{\par\vspace{8bp}}
\newcommand{\smallbody}{\fontsize{14bp}{21bp}\selectfont}
\newenvironment{datatable}{\begingroup\fontsize{16bp}{24bp}\selectfont\renewcommand{\arraystretch}{1.18}\setlength{\tabcolsep}{6bp}}{\endgroup}
\newcommand{\headerline}{\arrayrulecolor{SeedAccent}\hline\arrayrulecolor{SeedRule}}
\lstset{basicstyle=\ttfamily\fontsize{14bp}{21bp}\selectfont,columns=fullflexible,keepspaces=true,showstringspaces=false,breaklines=true,frame=none,aboveskip=4bp,belowskip=4bp}
\title{Nescity FPGAの高速化}
\author{あすた}
\filename{MICSCUP\_2026\_Asta}
\hypersetup{pdftitle={Nescity FPGAの高速化 5分版},pdfauthor={},pdfsubject={MAX10で計算サイクルを69.63倍改善したRTL評価},pdfkeywords={Nescity,FPGA,MAX10,MICS,229aed0}}

\begin{document}

% 01: 20 seconds. Use the template's original title layout.
\seedtitleframe[36]

% 02: 35 seconds
\begin{frame}{計算内容と課題}
\intent{隣のセルでも同じ近傍を読み直している}
\begin{columns}[T,onlytextwidth]
\begin{column}{0.60\textwidth}
32列 $\times$ 30行、合計960セル\par
各セルの次世代を5 $\times$ 5近傍から計算\par
\gap
上下左右は反対側につながる。\par
全セルは更新前の盤面を参照する。\par
\gap
\textbf{配布回路は24,000回の読出し}\par
$960\times25$ セルを毎世代読み直す。
\end{column}
\begin{column}{0.34\textwidth}
\textbf{近傍の重み}\par\vspace{4bp}
{\renewcommand{\arraystretch}{0.85}
$\begin{bmatrix}1&2&3&2&1\\2&3&4&3&2\\3&4&5&4&3\\2&3&4&3&2\\1&2&3&2&1\end{bmatrix}$}\par
\gap
$w=5-|dx|-|dy|$\par
重みの合計は65
\end{column}
\end{columns}
\end{frame}

% 03: 55 seconds
\begin{frame}{最適化(1/2)}
\intent{重なる近傍を使い回して、読出しを減らす}
\begin{datatable}
\begin{tabular}{@{}p{0.21\textwidth}p{0.49\textwidth}r@{}}
\textbf{段階} & \textbf{再利用するもの} & \textbf{読出し／世代}\\\headerline
配布回路 & セルごとに25近傍を集計 & 24,000\\\hline
列の再利用 & 5列のうち、重なる4列 & 5,400\\\hline
4行バッファ & 縦5セルのうち、過去4行 & 1,208\\\hline
行末の先読み & 行頭4列の集計 & \textcolor{SeedAccent}{\bfseries \FinalReads}\\\hline
\end{tabular}
\end{datatable}
\gap
追加の行バッファは\textbf{M9K 1個}、32ワード $\times$ 32 bit。\par
行末では保存した4列を使い、その間に次の行を先読みする。
\end{frame}

% 04: 55 seconds
\begin{frame}{最適化(2/2)}
\intent{別のセルの処理を重ねて、1クロックに1セル進める}
旧世代のバンクを読み、新世代のバンクへ同時に書く。\par
グリッドの容量は2 KiBのまま、1 KiBずつに分割する。
\gap
\begin{datatable}
\begin{tabular}{@{}lccccc@{}}
\textbf{計算クロック} & \textbf{SUM} & \textbf{AVG} & \textbf{BIAS} & \textbf{新値} & \textbf{書込み}\\\headerline
$n$   & セルE & セルD & セルC & セルB & \textcolor{SeedAccent}{\bfseries セルA}\\\hline
$n+1$ & セルF & セルE & セルD & セルC & \textcolor{SeedAccent}{\bfseries セルB}\\\hline
$n+2$ & セルG & セルF & セルE & セルD & \textcolor{SeedAccent}{\bfseries セルC}\\\hline
\end{tabular}
\end{datatable}
\gap
CPU側100 MHz・計算側10 MHzで、960セルを連続して書く。\par
完了通知は最終セルの書込み後に送る。
\end{frame}

% 05: 35 seconds
\begin{frame}{評価結果(1/2)}
\intent{計算クロックは10 MHzのまま、配布回路から約70倍}
\begin{datatable}
\begin{tabular}{@{}p{0.46\textwidth}r@{\hspace{16bp}}r@{}}
 & \textbf{配布回路} & \textbf{最適化後}\\\headerline
計算クロック数／世代 & \BaseCycles & \textcolor{SeedAccent}{\bfseries \FinalCycles}\\\hline
10 MHz換算 & \BaseMs\ ms & \textcolor{SeedAccent}{\bfseries \FinalMs\ ms}\\\hline
配布回路に対する速度 & 1.00倍 & \FinalSpeedup 倍\\\hline
\end{tabular}
\end{datatable}
\gap
計算サイクルを\CycleReduction\%削減。
\gap
測定区間はMMIOのSTARTからBUSY解除まで。\par
\textbf{描画やCRC処理を含む実機FPSは未測定。}
\end{frame}

% 06: 45 seconds
\begin{frame}{評価結果(2/2)}
\intent{QuartusでMAX10に収容できた。M9Kは残り2個}
\begin{datatable}
\begin{tabular}{@{}p{0.46\textwidth}r@{\hspace{16bp}}r@{}}
\textbf{NES全体の資源} & \textbf{使用数／総数} & \textbf{使用率}\\\headerline
LE & 6,020／8,064 & 75\%\\\hline
FF & 2,573／8,064 & 32\%\\\hline
M9K & \textcolor{SeedAccent}{\bfseries 40／42} & 95\%\\\hline
DSP（9-bit要素） & 0／48 & 0\%\\\hline
\end{tabular}
\end{datatable}
\gap
対象はMAX10 \texttt{10M08SAE144C8GES}。\par
アクセラレータのM9Kは\textbf{盤面2個＋行バッファ1個}。\par
全体の配置配線とSOF生成に成功した。
\end{frame}

% 07: 55 seconds
\begin{frame}{検証結果}
\intent{C版との一致に加え、実機でも0500までハッシュが一致}
\begin{itemize}\setlength{\itemsep}{4bp}
\item C版と\TestCases 配置・\TestGenerations 世代の全960セルで一致
\item 計算中の\ResetChecks 時点のリセットと、両バンクの照合を実施
\item 実機でも\textbf{0500（1,280世代）までハッシュ一致}を確認
\end{itemize}
\gap
{\smallbody
配置配線後のSetup slack：10 MHz側 $+4.040$ ns、100 MHz側 $-2.246$ ns。\par
100 MHz側にはタイミング違反が残る。}
\gap
\textbf{実装はこちら}\par
\href{https://github.com/hkaomua/MICS_CUP_2026-public}{\textcolor{SeedAccent}{\texttt{github.com/hkaomua/MICS\_CUP\_2026-public}}}\par
\end{frame}

\appendix
% A: optional
\begin{frame}{付録：最適化の経過}
\intent{読出しの再利用と処理の並行化を段階的に追加}
\begin{tikzpicture}
\begin{axis}[
width=0.80\textwidth,height=165bp,xbar,bar width=15bp,
xmin=0,xmax=82,xtick={0,20,40,60,80},
ytick={1,2,3,4,5},yticklabels={配布回路,列の再利用,先読みと演算の並行化,4行バッファ,連続パイプライン},
y dir=reverse,ymin=0.4,ymax=5.6,
axis x line*=bottom,axis y line*=left,axis line style={SeedRule},
xmajorgrids=true,grid style={SeedRule!40},tick style={draw=none},
xlabel={配布回路に対する速度倍率},
tick label style={font=\fontsize{12bp}{18bp}\selectfont},xlabel style={font=\fontsize{14bp}{21bp}\selectfont},
nodes near coords,point meta=explicit symbolic,every node near coord/.append style={font=\fontsize{14bp}{21bp}\selectfont,anchor=west,color=black},
]
\addplot[fill=SeedAccent,draw=none] coordinates {
(\StageARatio,1) [\StageALabel 倍]
(\StageBRatio,2) [\StageBLabel 倍]
(\StageCRatio,3) [\StageCLabel 倍]
(\StageDRatio,4) [\StageDLabel 倍]
(\StageERatio,5) [\StageELabel 倍]};
\end{axis}
\end{tikzpicture}
\gap
{\smallbody サイクル／世代：\StageACycles、\StageBCycles、\StageCCycles、\StageDCycles、\StageECycles。}
\end{frame}

% B: optional, proposal only
\begin{frame}{付録：2セル並列化の案}
\intent{残り2個のM9Kで約2倍を狙うが、ロジックが足りない可能性}
盤面を偶数行・奇数行に分け、2セルを同時に読み書きする。\par
読出し用・書込み用の盤面RAMを、合計2個から4個に増やす。
\gap
\begin{datatable}
\begin{tabular}{@{}p{0.15\textwidth}p{0.33\textwidth}l@{}}
\textbf{全体の資源} & \textbf{現在の使用量} & \textbf{2セル並列化の案}\\\headerline
M9K & 40／42個（95\%） & 42／42個、空きなし\\\hline
LE & 6,020／8,064（75\%） & 演算回路が増加、使用量は未算出\\\hline
LAB & 444／504（88\%） & 配置の余裕も少ない\\\hline
\end{tabular}
\end{datatable}
\gap
目標は同じ10 MHzで\textbf{約1.8～2倍}。未実装・未検証。\par
共通の演算を共有し、収容とタイミングが成立するかを確認する。
\end{frame}

\end{document}
